\chapter{把案情写成可检查的推导}
\label{chap:method}

本册按柏浪涛《2027精讲卷·刑法》的二十五讲组织。正文逐节展开，先给定义和条件，
再将题设事实代入，最后改变一个条件检验边界。每讲的记忆训练可独立闭卷使用。
附录保留较细的规则与边界供查阅。原书的教材观点、法条、司法解释和本册为演示而设的假定，分别标明。
覆盖范围是原书全部讲次、全部有标题的节及各节主要考点，个别罪名的全部量刑档、原书全部例题与脚注仍需结合教材和具体法源查阅。

\section{论域、时间和结论层次}

以一道题为单位建立记录
\[
 d=(c,p,a,o,e,t,v,j,\tau).
\]
$c$ 是题设案件，$p$ 是待评价的自然人或单位，$a$ 是需要单独判断的行为，$o$ 是行为对象，
$e$ 是相关结果，$t$ 是该规则要求的时点，$v$ 是规范版本，$j$ 是法域或地域参数，
$\tau$ 是明确选用的解释方案。同一案件有数名行为人、数个行为或数个结果时，分别建立记录。
例如 $\mathrm{Age}(p,t)$ 输出周岁数，$\mathrm{Cause}(a,e)$ 是真假命题。
$\mathrm{Age}(p,t)\ge16$ 可以作为一个年龄条件，不能拿它代替 $\mathrm{Cause}(a,e)$。

式中 $\land,\lor,\neg$ 分别表示合取、析取、否定。$A\lor B$ 允许两项同真，只有明确互斥时再加 $\neg(A\land B)$。
有限合取 $\bigwedge_{i=1}^m A_i$ 表示每一项都要审查。$\exists x$ 需要提供见证对象，$\forall x$ 须界定对象范围。
$A:=B$ 是记号定义，$A\Rightarrow B$ 是一个方向的条件命题，$A\Leftrightarrow B$ 须有两个方向的依据。
$\Gamma\vdash_{\mathcal R} A$ 记录从已列前提集 $\Gamma$ 使用规则集 $\mathcal R$ 得到 $A$ 的推导。
不同章节可另定义符号，同一符号的含义以其所在小节为准。

一份答题记录要分别填入五类输出：
\begin{enumerate}
 \item 对特定罪名的\textbf{构成符合性}，包括主体、行为、结果、因果和相应罪过。
 \item \textbf{违法性与个人责任}，包括正当防卫、责任年龄、责任能力等审查。
 \item \textbf{形态和参与范围}，包括预备、未遂、中止、既遂及各人的参与范围。
 \item \textbf{最终评价}，包括多个候选罪名之间的竞合、特别规定及罪数处理。
 \item \textbf{法律后果}，包括法定刑档、裁量授权、执行或追诉的条件。
\end{enumerate}
采用教材的改良版四要件分析时，可将这五类结果作为答题顺序。
这是一份检查清单，遇到特别条款或当事人身份差异时，仍须重新写出具体的依赖关系。

\begin{center}
\begin{tikzpicture}[
  node distance=6mm and 5mm,
  logicstep/.style={draw=GuideGreen,rounded corners=1.5pt,align=center,
    minimum height=8mm,text width=19mm,inner sep=2pt,font=\small},
  relation/.style={-{Stealth[length=2mm]},draw=GuideGreen,semithick}]
 \node[logicstep] (f) {固定人、事\\时间、版本};
 \node[logicstep,right=of f] (k) {逐项核对\\候选构成};
 \node[logicstep,right=of k] (r) {审查违法\\与个人责任};
 \node[logicstep,right=of r] (s) {确定形态\\参与与罪数};
 \node[logicstep,right=of s] (p) {列明处罚\\与裁量范围};
 \draw[relation] (f) -- (k);
 \draw[relation] (k) -- (r);
 \draw[relation] (r) -- (s);
 \draw[relation] (s) -- (p);
\end{tikzpicture}
\end{center}

\section{证据状态和命题真假}

先登记待证明的事实。对命题 $A$，分别记录肯定证据是否达到当前用途的支持要求和否定证据是否达到该要求。
用二元状态 $\sigma(A)=(s_+,s_-)$ 表示这个登记结果：
\[
\begin{array}{c|c|l}
 \sigma(A)&\text{记号}&\text{填表含义}\\\hline
 (1,0)&\mathsf T&\text{有可采的肯定根据，未见同级冲突}\\
 (0,1)&\mathsf F&\text{有可采的否定根据，未见同级冲突}\\
 (0,0)&\mathsf U&\text{材料不足，肯定与否定均未确立}\\
 (1,1)&\mathsf X&\text{支持材料发生尚待解决的冲突}
\end{array}
\]
对 $\mathsf U$，应写明缺哪项事实。对 $\mathsf X$，应查明证据或者暂列两套相容的题设继续条件推导。
只把已确认的命题及否定命题放入一致的 $\Gamma$，以免把冲突材料直接输入经典逻辑。
例如“未能证明乙事先通谋”最多表示肯定命题的证明未完成。若乙在行为结束后才第一次得知事实，才有独立根据排除事先通谋。

刑事定罪证据标准依据《刑事诉讼法》第55条。需要证据证明定罪量刑的事实，
核实据以定案的证据，并在综合全案后排除合理怀疑。
$\sigma$ 只是学习时登记推理输入的方法。遇到有利于被追诉人的材料，是否满足证明标准和如何处理疑问，
按刑事诉讼法及相应规范审查。不能借“未证明例外”把依法应调查、应排除的疑问转给被追诉人证明。
本册推导例通常直接给定事实，以便把证明事实和适用规则两个任务分开。

\section{一条规范如何进入一份证明}

把一条可适用规则写成
\[
 r=(\mathrm{id},\text{法源或解释},v,j,I,\tau,
       A_r,B_r,H_r).
\]
$I$ 是相应的时间范围，$A_r$ 是积极条件，$B_r$ 是阻却或者例外条件，$H_r$ 是在该路径下的后果。
审查后，将适用性记为 $\mathrm{Scope}_r(d)$。对于前提和例外均已确定、且没有未处理规则冲突的例题，
可按下列推理形式使用 $r$：
\[
\frac{
 \Gamma\vdash \mathrm{Scope}_r(d)
 \qquad \Gamma\vdash A_r(d)
 \qquad \Gamma\vdash \neg B_r(d)
}{
 \Gamma\vdash_{\mathcal R} H_r(d)
}\quad(r).
\]
若规则含有多个备选分支，先选一条有事实支持的分支。若有特别法条、司法解释或者并罚规则，明示它作用于哪一层。
$H_r$ 可以是“构成条件具备”“应当减轻”或“可以宣告缓刑”。这些是不同类型的输出。
若 $\mathrm{May}(X)$ 表示可依法选 $X$，则 $\mathrm{May}(X)\nRightarrow X$。
例如缓刑的积极条件已具备，仍需核对法定禁止项，之后才判断是裁量授权还是特殊主体的应当宣告规则。

\paragraph{推导为何成立。}
固定同一法源版本和解释方案，把每条已审查可用的规则解释成在该范围内保真的推理步骤。
考虑一棵有限证明树。叶节点是给定或已确认的事实，真值已固定。若某个内部节点的子节点都成立，
则其所标的合取消去、析取引入或条件推理保证该节点成立。按树的高度归纳，根节点随之成立。
因此结论依赖于每个叶节点的可靠性和每条规则的适用性。换掉一个事实或规范版本，须从受影响的节点重新推演。
这个证明说明形式推导的条件正确性，实际规则是否适用仍须依法解释、查证。

\paragraph{必要条件与充分条件。}
若 $H\Rightarrow A$，$A$ 是 $H$ 的必要条件。为检验逆命题，找一组使 $A$ 真而 $H$ 假的相容事实。
若 $H$ 需 $A\land B$，赋值 $A=1,B=0$ 就足以否定“仅有 $A$ 即有 $H$”。
同理，从 $A\Rightarrow B$ 和 $\neg A$ 不能推出 $\neg B$。取 $A=0,B=1$，前提真而所猜的结论假。
这是分则中“未命中这一罪名分支，继续查其他分支”的逻辑原因。

\paragraph{学说分歧如何并行。}
保留共同事实 $\Gamma$，为两种观点各列规则 $\mathcal R_\tau,\mathcal R_{\tau'}$。
如果只在 $\mathcal R_\tau$ 中能够推出 $H$，写成“依 $\tau$ 可得 $H$”，并指出分歧在哪个条件。
若所有题设承认的解释都推出 $H$，可以把它记为本题在这些解释下的共同结论。
两套规则不能互相借用前提以拼成第三套结论。可用同一题设对照第八讲的认识错误和第十二讲的共犯处理。

\section{示范证明：防卫过当为什么要核对两道门槛}

设 $B$ 为一般防卫的起因、时间、对象、目的等基础条件成立。$S$ 为命中特殊防卫条款。
$L$ 为防卫行为明显超过必要限度，$D$ 为造成重大损害。
在符合基础条件且明确不进入特殊防卫的模型中，按《刑法》第20条第二款及正当防卫指导意见，将过当条件记为
\[
 B\land\neg S\quad\text{固定时},\qquad O:=L\land D.
\]
已固定范围中的四种组合如下：
\[
\begin{array}{cc|c|l}
 L&D&O&\text{判断}\\\hline
 0&0&0&\text{未命中}\\
 0&1&0&\text{损害重大，但限度条件欠缺}\\
 1&0&0&\text{明显超限，但重大损害条件欠缺}\\
 1&1&1&\text{进入防卫过当及相应责任审查}
\end{array}
\]
假设题目给定正在发生的非法侵害，且甲为制止侵害而针对侵害人实施防卫。题设另明确不是严重危及人身安全的暴力犯罪，
并经具体情境评价认定防卫手段并未明显超限，即 $\Gamma=\{B,\neg S,\neg L,D\}$。
由 $O\Rightarrow L$ 和 $\neg L$ 可推出 $\neg O$。因此“造成重大损害”本身不推出防卫过当。
若把唯一变项改为 $L=1$，合取引入给出 $L\land D$，再依据定义推出 $O$。
之后仍须审查相应罪名的罪过、因果、责任；依法负刑事责任时，适用该款的应当减轻或者免除处罚。
若改为侵害已经结束，则应回到 $B$，不能在假定 $B$ 成立的表格中继续代入。

\section{把推导变成闭卷记忆}

\paragraph{先写提取卡，再背口诀。}
每节准备一张卡。正面写题设触发事实和一个易混对象，背面固定五栏：\textbf{目标结论、必要前提、并列分支、例外、一个反例}。
例如防卫过当卡正面写“造成重大损害，必然过当吗？”，背面依次填“先查防卫基础和特殊防卫”“明显超限”“重大损害”“两项同真才命中”“重大损害而未明显超限”。
短句“先查防卫，再验双门”只用来定位检查顺序，闭卷时要能把两个门完整展开。

\paragraph{在三分钟内重建规则。}
第一遍看正文，遮住公式，用中文列出各前提。第二遍只看卡片正面，写公式并标记“且”或“或”。
第三遍用一句“事实 $\to$ 条件 $\to$ 结论”解释每个推理节点。写完换一个事实重做，
如改动时间点、年龄、是否有处分意识或是否事前通谋。仍得相同结论时，说明哪个分支继续支撑它。

\paragraph{按错误类型复习。}
可在当日、第1日、第3日、第7日、第14日闭卷复测。每次先写再翻答案。
只抄自己漏掉的条件和正确方向，按“漏条件”“混淆层次”“误用时间”“把未知当否定”“把裁量当必然”分别记错。
下一次先做上一轮错项，再做一组近邻对比。成功标准是能在没有选项提示时写全前提，指出最先改变的推理节点，并说出该节点改变后的后续审查。

\paragraph{一张卡的迁移练习。}
写出“只有达到规定年龄才可能对某罪负刑责”的方向。把题中年龄从十五周岁改为十六周岁，
你只应更新年龄入口。对象、行为、罪过、阻却和责任能力的审查仍各有对应事实。
若年龄本身不明，就把 $\mathrm{Age}(p,t)$ 列入待证栏，不能补写成十五岁或十六岁。
